\documentclass[11pt]{article}

\usepackage[margin=1in]{geometry}
\usepackage{amsmath}
\usepackage{amssymb}
\usepackage{bm}
\usepackage{hyperref}

\title{DynBBH Tabulated Trajectories, Analytic Backgrounds, and Smooth Excision}
\author{}
\date{\today}

\begin{document}
\maketitle

\section{Scope}

This note summarizes the changes made to the boosted Kerr--Schild
binary-black-hole problem generator and to the associated excision path.
The historical circular Keplerian binary remains the default.  New behavior
is enabled only by explicit runtime flags, chiefly the trajectory-table flag,
smooth-excision flag, and optional magnetic-damping flag.

The main additions are:
\begin{itemize}
  \item a tabulated black-hole trajectory/property path with masses,
  positions, spins, and coordinate velocities;
  \item cubic-Hermite interpolation of positions and velocities, with
  analytic time derivatives used by the metric and extrinsic-curvature
  calculation;
  \item refactored boosted Kerr--Schild algebra built from reusable boost,
  Kerr--Schild, and superposition helpers;
  \item a smooth hydrodynamic drain inside the apparent horizons;
  \item an optional constraint-transport-compatible magnetic damping term
  applied through edge-centered EMFs.
\end{itemize}

The smooth hydrodynamic drain acts on primitive and conserved matter
variables.  It does not directly rescale, zero, or window the face-centered
magnetic field.  Direct pointwise modification of the staggered magnetic
field can create discrete magnetic monopoles, so magnetic damping is instead
implemented as an EMF contribution to the constrained-transport update.

\section{Trajectory Tables and Analytic Derivatives}

When the trajectory-table mode is disabled, the binary state is still the
historical circular Keplerian orbit.  When it is enabled, the table columns
are interpreted as
\begin{equation}
\begin{split}
&t,\; m_1,\; m_2,\; X_1,\;Y_1,\;Z_1,\;X_2,\;Y_2,\;Z_2,\;
\chi_{1x},\;\chi_{1y},\;\chi_{1z},\\
&\chi_{2x},\;\chi_{2y},\;\chi_{2z},\;
V_{1x},\;V_{1y},\;V_{1z},\;V_{2x},\;V_{2y},\;V_{2z}.
\end{split}
\end{equation}
The table spin entries are dimensionless spins \(\chi_A^i\).  The Kerr
spin vectors used by the metric and horizon shape are
\begin{equation}
  a_A^i = m_A\chi_A^i .
\end{equation}

For each position component \(x(t)\) on a segment
\([t_0,t_1]\), let \(h=t_1-t_0\) and
\(s=(t-t_0)/h\).  The interpolant is the cubic Hermite polynomial
\begin{equation}
\begin{split}
  x(t) ={}&
  (2s^3-3s^2+1)x_0
  +(s^3-2s^2+s)h v_0\\
  &+(-2s^3+3s^2)x_1
  +(s^3-s^2)h v_1 .
\end{split}
\end{equation}
The first derivative gives the interpolated velocity and the second
derivative gives the acceleration used by the analytic time derivatives.
Masses and spin components are linearly interpolated, with piecewise-constant
time derivatives on each segment:
\begin{equation}
  f(t) = (1-s)f_0 + s f_1,
  \qquad
  \dot f(t) = \frac{f_1-f_0}{h}.
\end{equation}
The active table segment is cached so repeated calls at nearby times do not
perform a full search.

The metric path now evaluates the four-metric and its derivatives
analytically through dual numbers rather than by finite-differencing metric
evaluations in time.  In particular, the boost Jacobian has no zero-velocity
early return.  Instead, the regular small-\(V^2\) expansion
\begin{equation}
  \frac{\gamma - 1}{V^2}
  =
  \frac{1}{2}
  +\frac{3}{8}V^2
  +\frac{5}{16}V^4
  +O(V^6)
\end{equation}
is used.  This preserves dual-number derivatives of the boost even when an
instantaneously stationary black hole has nonzero acceleration.

\section{Boosted Kerr--Schild Superposition}

Each black-hole contribution is assembled as a boosted Kerr--Schild
perturbation and then superposed on the background.  The algebra is factored
into helper routines for the Lorentz boost, spatial Kerr--Schild radius,
single-hole perturbation, and final ADM conversion.  This avoids duplicating
large expanded symbolic expressions and makes the table-driven and analytic
Keplerian paths share the same metric assembly.

For GPU safety, table states used inside Kokkos kernels are copied into
plain by-value state objects before launching the kernel.  Device lambdas
therefore capture POD data and Kokkos views, not host stack references.

\section{Boosted Spinning Horizon Mask}

For black hole \(A\), let \(X_A^i\), \(V_A^i\), \(\chi_A^i\), and \(m_A\)
denote the position, coordinate velocity, dimensionless spin vector, and
mass.  The Kerr spin vector used in the horizon and metric algebra is
\begin{equation}
  a_A^i = m_A\chi_A^i .
\end{equation}
For a cell center \(x^i\), define the displacement
\begin{equation}
  \Delta x_A^i = x^i - X_A^i .
\end{equation}
The displacement is mapped into the instantaneously boosted local frame by
\begin{equation}
  \Delta \bar{x}_A^i
  =
  \Delta x_A^i
  +
  \frac{\gamma_A - 1}{V_A^2}
  V_A^i
  \left(V_{A j}\Delta x_A^j\right),
  \qquad
  \gamma_A = \frac{1}{\sqrt{1 - V_A^2}} .
\end{equation}
For small \(V_A^2\), the regular expansion
\begin{equation}
  \frac{\gamma_A - 1}{V_A^2}
  =
  \frac{1}{2}
  + \frac{3}{8} V_A^2
  + \frac{5}{16} V_A^4
  + O(V_A^6)
\end{equation}
is used to avoid loss of precision.

The spinning Kerr--Schild radial coordinate is computed from
\begin{equation}
  r_A^2 =
  \frac{
    \bar{r}_A^2 - a_A^2
    +
    \sqrt{\left(\bar{r}_A^2-a_A^2\right)^2 + 4a_A^2 z_A^2}
  }{2},
\end{equation}
where
\begin{equation}
  \bar{r}_A^2 = \Delta \bar{x}_{A i}\Delta \bar{x}_A^i,
  \qquad
  z_A = \frac{a_{A i}\Delta \bar{x}_A^i}{|a_A|}.
\end{equation}
If \(a_A^2=0\), this reduces to the Euclidean radius
\(\sqrt{\Delta \bar{x}_{A i}\Delta \bar{x}_A^i}\).

The isolated Kerr horizon radius is approximated by
\begin{equation}
  r_{H,A} = m_A + \sqrt{m_A^2 - |a_A|^2}.
\end{equation}
A cell is marked as geometrically excised when
\begin{equation}
  \min_A \left(r_A - r_{H,A}\right) \leq 0 .
\end{equation}

\section{Smooth Excision Weight}

The smooth drain weight for hole \(A\) is
\begin{equation}
  W_A = S\!\left(\frac{r_{H,A}-r_A}{\Delta r}\right),
\end{equation}
where \(\Delta r\) is the configured smooth-excision width.  If no width is
given in the input file, the code uses a conservative default
\begin{equation}
  \Delta r_{\rm default}
  =
  \max\!\left(0.25\,r_{\rm exc},\,2\Delta x_{\max}\right),
\end{equation}
with \(\Delta x_{\max}\) the largest active mesh spacing across all MPI
ranks.  The local block-pack maximum is reduced with a global maximum so the
default is rank independent on decomposed AMR or uneven MPI layouts.
The global cell weight is
\begin{equation}
  W = \max_A W_A .
\end{equation}
The limiter \(S\) is the cubic smoothstep
\begin{equation}
  S(q) =
  \begin{cases}
    0, & q \leq 0,\\
    3q^2 - 2q^3, & 0 < q < 1,\\
    1, & q \geq 1.
  \end{cases}
\end{equation}
Thus \(W=0\) outside the horizon and \(W\to 1\) deeper inside the excision
region.

For lapse-based excision, the same construction is applied in lapse space:
\begin{equation}
  W = S\!\left(\frac{\alpha_{\rm exc} - \alpha}{\Delta \alpha}\right),
\end{equation}
where \(\alpha_{\rm exc}\) is the excision lapse threshold and
\(\Delta \alpha\) is the configured lapse width.

The floor mask and conservative-flux mask are deliberately distinct.  For
puncture excision, the floor mask uses \(r_A \le r_{H,A}\), while the flux
mask uses
\begin{equation}
  r_A \le \max(r_{H,A}, r_{\rm flux}).
\end{equation}
For lapse excision with smooth excision enabled, the flux mask is buffered to
\begin{equation}
  \alpha < \alpha_{\rm exc} + \Delta\alpha .
\end{equation}
Without smooth excision, the lapse flux mask reduces to the historical
\(\alpha < \alpha_{\rm exc}\) mask.

\section{Primitive Variable Drain}

Let \(A_\ast\) be the nearest black hole center to the cell.  The target
velocity is the black-hole coordinate velocity plus an inward drain speed:
\begin{equation}
  v_{\rm targ}^i
  =
  V_{A_\ast}^i
  -
  v_{\rm in}
  \frac{x^i - X_{A_\ast}^i}
       {\sqrt{(x^j-X_{A_\ast}^j)(x_j-X_{A_\ast j})}} .
\end{equation}
This is a coordinate three-velocity.  The GR hydro and GRMHD primitive
variables store the Wilson/Valencia spatial four-velocity \(Wv^i\), so the
target is converted before blending:
\begin{align}
  v_{\rm targ}^2 &= \gamma_{ij}v_{\rm targ}^i v_{\rm targ}^j,\\
  v_{\rm targ}^i &\leftarrow
  v_{\rm targ}^i
  \min\!\left(1,\sqrt{\frac{v_{\max}^2}{v_{\rm targ}^2}}\right),\\
  W_{\rm targ} &= \frac{1}{\sqrt{1-v_{\rm targ}^2}},\\
  u_{\rm targ}^i &= W_{\rm targ}v_{\rm targ}^i .
\end{align}
Here \(v_{\max}^2 < 1\) is chosen consistently with the local velocity
ceiling when one is available, and otherwise with a tight subluminal cap.
The primitive variables are blended toward an excision state,
\begin{align}
  \rho' &= (1-W)\rho + W\rho_{\rm targ},\\
  u'^i &= (1-W)u^i + Wu_{\rm targ}^i,\\
  p' &= (1-W)p + Wp_{\rm exc}.
\end{align}
When no magnetization cap is requested,
\begin{equation}
  \rho_{\rm targ} = \rho_{\rm exc}.
\end{equation}

\section{Optional Magnetization Cap}

The smooth excision can optionally cap the cell-centered proxy
\(\sigma \equiv B_{\rm cc}^2/\rho\).  This is controlled by
\(\sigma_{\max}\).  If \(\sigma_{\max}>0\),
\begin{equation}
  \rho_{\rm targ}
  =
  \max\!\left(
    \rho_{\rm exc},
    \frac{B_{\rm cc}^2}{\sigma_{\max}}
  \right).
\end{equation}
If \(\sigma_{\max}<0\), the cap is disabled.  The value
\(\sigma_{\max}=0\) is rejected.

The magnetic field entering the cap is the existing cell-centered average of
the face-centered field,
\begin{equation}
  B_{\rm cc}^i =
  \frac{1}{2}\left(B^i_{f,-} + B^i_{f,+}\right).
\end{equation}
The cap therefore changes the matter state, not the staggered magnetic
field.  After the primitive variables are blended, the conserved fluid
variables are recomputed using the local metric and the unchanged magnetic
field:
\begin{equation}
  U'_{\rm hydro} =
  {\rm PrimToCon}\!\left(P', B_{\rm cc}, g_{\mu\nu}\right).
\end{equation}

\section{Magnetic Constraint Preservation}

The face-centered magnetic field is not directly modified by the excision
operator:
\begin{equation}
  B_f' = B_f .
\end{equation}
All changes to \(B_f\) remain those produced by the constrained-transport
update.  Schematically, the CT update is
\begin{equation}
  B_f^{n+1}
  =
  B_f^n
  -
  \Delta t\, C_h(E_e),
\end{equation}
where \(C_h\) is the discrete curl mapping edge-centered EMFs \(E_e\) to
face-centered magnetic fluxes.  Applying the discrete divergence operator
\(D_h\) gives
\begin{equation}
  D_h B_f^{n+1}
  =
  D_h B_f^n
  -
  \Delta t\,D_h C_h(E_e).
\end{equation}
On the staggered CT mesh,
\begin{equation}
  D_h C_h(E_e) = 0
\end{equation}
as a discrete identity.  Therefore
\begin{equation}
  D_h B_f^{n+1} = D_h B_f^n .
\end{equation}

Because the smooth excision operator never multiplies or zeros the
face-centered magnetic field, it does not introduce a new source term in the
discrete magnetic divergence equation.  The optional magnetization cap also
preserves this property: it modifies \(\rho_{\rm targ}\), not \(B_f\).

\section{CT-Compatible Magnetic Damping Design}

Magnetic-field damping inside the horizon should be implemented through an
edge-centered EMF contribution, not through a pointwise face-field source.
The minimal CT-compatible form is an additional EMF
\begin{equation}
  E_e = E_{e,{\rm ideal}} + E_{e,{\rm damp}},
\end{equation}
inserted after the ideal corner EMFs are built and before EMF exchange,
restriction, and CT.  A resistive damping choice is
\begin{equation}
  E_{e,{\rm damp}} = \eta_e W_e J_e,
  \qquad
  J = \nabla_h \times B_f ,
\end{equation}
where \(W_e\) is the smooth excision weight interpolated to the edge,
\(\eta_e \ge 0\) is a damping coefficient, and \(J_e\) is a compatible
edge-centered current built from the staggered face field.  The magnetic
update remains
\begin{equation}
  B_f^{n+1}
  =
  B_f^n
  -
  \Delta t\, C_h(E_{e,{\rm ideal}} + E_{e,{\rm damp}}),
\end{equation}
so
\begin{equation}
  D_h B_f^{n+1}
  =
  D_h B_f^n
  -
  \Delta t\,D_h C_h(E_{e,{\rm ideal}} + E_{e,{\rm damp}})
  =
  D_h B_f^n .
\end{equation}
The discrete divergence remains unchanged because the added term still
enters only through the discrete curl.

For adaptive mesh refinement, the damping EMF must be part of the same
edge-centered EMF arrays that are exchanged and corrected across coarse--fine
interfaces.  The intended task order is therefore
\begin{equation}
  {\rm CornerE}
  \rightarrow
  {\rm AddUserDampingE}
  \rightarrow
  {\rm SendE/RecvE/CorrectE}
  \rightarrow
  {\rm CT}.
\end{equation}
The damping coefficient should obey the explicit diffusion stability limit,
schematically
\begin{equation}
  \eta_e \lesssim C_{\rm diff}\frac{\Delta x_e^2}{\Delta t}.
\end{equation}
The implemented runtime controls are
\begin{align}
  {\tt coord/smooth\_excision\_b\_damping} &\in \{{\tt true},{\tt false}\},\\
  {\tt coord/smooth\_excision\_b\_damping\_eta} &= \eta_0,\\
  {\tt coord/smooth\_excision\_b\_damping\_cfl} &= C_{\rm diff}.
\end{align}
When \(C_{\rm diff}>0\), the active coefficient is capped as
\begin{equation}
  \eta_e =
  \min\!\left(\eta_0,\,
  C_{\rm diff}\frac{\Delta x_{\min}^2}{\Delta t}\right).
\end{equation}
Setting \(C_{\rm diff}=0\) disables the cap.
This resistive term damps current-carrying structure.  If a true flux drain
is later desired, an artificial advective EMF
\begin{equation}
  E_{e,{\rm adv}} = -v_{B,e}\times B_e
\end{equation}
can also be added at edges, with \(v_{B,e}\) directed inward inside the
horizon.  It is also CT-compatible, but it transports magnetic flux rather
than deleting it locally.

\section{Robustness and Safety Fixes}

Several implementation details were tightened during review:
\begin{itemize}
  \item The trajectory state captured by initialization kernels is a
  by-value POD object, not a reference to host stack storage.
  \item The boost Jacobian uses the small-\(V^2\) series above instead of
  returning an identity matrix at exactly zero velocity, so analytic time
  derivatives are retained.
  \item The Valencia cooling source now returns immediately unless both MHD
  and ADM storage are present, avoiding a hydro-only null dereference.
  \item The vector-potential transformation guards the cylindrical
  denominator \(x^2+y^2\) with a positive floor before forming
  \(A_x\) and \(A_y\), preventing \(0/0\) on the rotation axis.
  \item Smooth-excision drain targets are converted from coordinate velocity
  to the primitive Wilson/Valencia variable \(Wv^i\) before blending and are
  capped with the configured velocity ceiling.
  \item The post-blend primitive state is checked for positivity and
  finiteness over the active primitive/scalar slots before it is accepted.
  \item The optional user electric-field callback is dispatched only when
  both its enable flag and function pointer are set.
  \item The smooth-excision width default is made MPI consistent by a global
  maximum over the local mesh spacings.
\end{itemize}

\section{Unresolved-Horizon Sink Fallback}

If the local mesh spacing around a hole is too large compared with that
hole's horizon diameter, the puncture excision mask is formally
under-resolved.  With the \texttt{problem/unresolved\_sink} flag, the problem
generator keeps puncture excision enabled and applies an additional
conservative sink over a larger radius.  This path is intended for cheap
early inspiral evolutions where one or both holes are deliberately unresolved
before a later AMR restart resolves the horizons.

The decision is made separately for each hole.  Let
\(\Delta x_A\) be the largest cell spacing in the active MeshBlock containing
hole \(A\).  The hole is treated as resolved when
\begin{equation}
  \frac{2r_{H,A}}{\Delta x_A} \ge N_{\rm H},
  \qquad N_{\rm H}=20 \quad \hbox{by default}.
\end{equation}
When this condition is satisfied, the sink for that hole is disabled.  The
puncture excision radius is never suppressed by the sink path: excision stays
active whenever \texttt{coord/bh\_excise=true}.  If only one hole is resolved,
only that hole's sink is disabled; the other hole keeps its sink while both
holes continue to use the configured excision masks.

For hole \(A\), define the smooth sink weight
\begin{equation}
  W_A^{\rm sink}
  =
  S\!\left(\frac{R_{\rm sink}-|x-X_A|}{\Delta R_{\rm sink}}\right),
  \qquad
  W^{\rm sink} = \max_A W_A^{\rm sink},
\end{equation}
with the same cubic smoothstep \(S\) used by smooth excision.  The automatic
sink radius is
\begin{equation}
  R_{{\rm sink},A}
  =
  \max\!\left(
    N_R \Delta x_A,\;
    R_{\rm sink,min},\;
    r_{H,A}
  \right),
  \qquad N_R=10 \quad \hbox{by default},
\end{equation}
so the active sink is at least 20 cells wide, or 10 cells in radius.  As AMR
refines the hole, \(\Delta x_A\) decreases and the sink radius shrinks
automatically.  The optional configured \texttt{sink\_radius} acts as a
minimum physical radius; setting it to zero gives a purely resolution-scaled
sink.  Over a step \(\Delta t\), the damping fraction is
\begin{equation}
  f_{\rm sink}
  =
  1-\exp\!\left(-\frac{\Delta t\,W^{\rm sink}}{\tau_{\rm sink}}\right).
\end{equation}
The conservative matter variables are damped as
\begin{align}
  D' &= (1-f_{\rm sink})D + f_{\rm sink}D_{\rm floor},\\
  S'_i &= (1-f_{\rm sink})S_i,\\
  \tau' &= (1-f_{\rm sink})\tau + f_{\rm sink}\tau_{\rm floor},\\
  (\rho Y_s)' &= (1-f_{\rm sink})(\rho Y_s),
\end{align}
where \(D_{\rm floor}\) and \(\tau_{\rm floor}\) are configured sink targets
or fall back to the EOS floors.  The sink does not modify the face-centered
magnetic field.  Avoiding a primitive-to-conserved conversion through the
unresolved puncture metric is important: near \(r=0\), such a conversion can
inject an artificially large conserved energy even for a small target
pressure.

This fallback is not a replacement for horizon excision.  It is a coarse-grid
drain that removes matter and momentum smoothly while preserving the
constrained-transport representation of \(B_f\).  On a restarted zoom run, the
same per-hole logic continues to apply after AMR regridding: the sink radius
follows the local resolution and each sink turns off only after that hole
satisfies the configured cells-across-horizon criterion.

\section{Device, MPI, and AMR Safety}

The new Kokkos kernels capture only scalars, POD state, and Kokkos views.
No STL containers, host arrays, or host-only pointers are dereferenced inside
device lambdas.  The tabulated trajectory itself remains host-side; only the
interpolated state needed by a kernel launch is copied into a device-safe
value object.

The magnetic damping is MPI and AMR compatible because it modifies the same
edge-centered EMF arrays used by the normal CT update.  The task ordering is
\begin{equation}
  {\rm CornerE}
  \rightarrow
  {\rm EFieldSrc}
  \rightarrow
  {\rm SendE/RecvE/CorrectE}
  \rightarrow
  {\rm CT},
\end{equation}
so damped EMFs are exchanged across MPI boundaries and corrected across
coarse--fine interfaces before they update \(B_f\).

\section{Validation Performed}

The changes were checked with serial and MPI builds of the dynbbh problem
generator.  A two-rank smooth-excision and magnetic-damping run completed
with finite primitive and magnetic fields.  The inspected constrained
transport diagnostic remained at roundoff,
\begin{equation}
  \max |D_h B_f| \simeq 1.6\times 10^{-16}.
\end{equation}
A separate discrete CT damping-rate check matched the analytic one-step
diffusion amplification to roundoff while preserving zero discrete
divergence.  Long stronger-field runs showed the prescribed damping reduces
the magnetic-field strength inside the excision region relative to the
undamped run while preserving the CT divergence constraint.

\section{What Is Not Included}

No direct magnetic window function is included in this implementation.  The
optional magnetic damping path modifies only edge-centered EMFs and is disabled
by default.

\end{document}
